\section{Modelo orientado a Accuracy}

El diseño inicial de esta fase se basó en una cadena \textit{RFE + MLP}: primero se realizó selección de variables con RFE (estimador lineal) y, sobre ese subconjunto, se entrenó una red neuronal multicapa optimizada con validación cruzada estratificada. Como paso posterior, y manteniendo el mismo protocolo experimental, se evaluaron modelos alternativos de ensamble para intentar mejorar la métrica objetivo (\textit{Accuracy}).

El criterio de selección se mantuvo constante durante toda la fase: elección por rendimiento en validación interna y evaluación final en el conjunto de test bloqueado, sin recalibrar umbrales con información de test.

\subsection{Selección de variables con RFE}

Sobre la representación preprocesada del bloque lineal se aplicó RFE con regresión logística, fijando $k=5$ variables. El subconjunto seleccionado en esta etapa fue:

\begin{itemize}
	\item \texttt{nominal\_\_V1\_4.0}
	\item \texttt{nominal\_\_V1\_24.0}
	\item \texttt{nominal\_\_V1\_31.0}
	\item \texttt{numeric\_\_V20}
	\item \texttt{numeric\_\_V8}
\end{itemize}

Este bloque produjo un modelo base competitivo y permitió establecer una referencia inicial para Accuracy en test.

\subsection{Ajuste del modelo y selección final}

En una primera iteración se optimizó una MLP con \texttt{GridSearchCV} sobre arquitectura, activación, regularización y tasa de aprendizaje. Esa variante obtuvo un rendimiento de referencia de aproximadamente \textbf{0.8559} en Accuracy sobre test.

Dado que existía margen de mejora, se ejecutó una segunda iteración de búsqueda sobre modelos de ensamble de árboles (Random Forest y Extra Trees), manteniendo validación cruzada estratificada y selección de umbral en validación interna. El mejor candidato por validación cruzada fue \textbf{Random Forest}, adoptado como modelo final de Accuracy.

\begin{table}[h]
\centering
\caption{Resultados finales del modelo orientado a Accuracy (test)}
\begin{tabular}{|l|c|}
\hline
Modelo final & \texttt{RF\_Accuracy} \\
\hline
Umbral de decisión & 0.49 \\
\hline
Accuracy & 0.9492 \\
\hline
Recall & 0.8156 \\
\hline
Precision & 0.9200 \\
\hline
F1-score & 0.8647 \\
\hline
ROC-AUC & 0.9211 \\
\hline
PR-AUC & 0.8968 \\
\hline
Matriz de confusión $(TN,FP,FN,TP)$ & $(557,10,26,115)$ \\
\hline
\end{tabular}
\end{table}

La mejora respecto al baseline MLP fue sustancial: $0.9492 - 0.8559 = 0.0933$ puntos absolutos de Accuracy. Este resultado confirma que, para el objetivo de calidad global de clasificación, el ensamble de árboles captura mejor la estructura no lineal del problema en comparación con el modelo neuronal inicial.
